%% Dualite en grille 2x2 : en haut, la porte Pi_theta (temps, variable
%% t) et sa TF theta*sinc(pi f theta) (frequence, variable f) ; en bas,
%% en repartant du sinus cardinal vu comme un signal TEMPOREL (variable
%% t), sa TF redonne une porte, maintenant vue en FREQUENCE (variable
%% f) -- Pi_theta est paire donc x(-f)=x(f). Meme paire de courbes,
%% roles temps/frequence inverses : l'axe des abscisses de chaque
%% sous-figure porte volontairement la bonne variable. Consomme par
%% slides.
\begin{tikzpicture}
    \begin{groupplot}[
        group style={group name=dualgrid, group size=2 by 2, horizontal sep=1.7cm, vertical sep=1.8cm},
        width=5.6cm, height=3.5cm,
        ytick=\empty,
        axis lines=middle,
    ]
    \nextgroupplot[xlabel={$t$}, xmin=-2.2,xmax=2.2, ymin=-0.2,ymax=1.5]
    \addplot[niceblue, thick] coordinates {(-2,0)(-0.5,0)(-0.5,1)(0.5,1)(0.5,0)(2,0)};
    \node at (axis cs: 0,1.7) {\small $\Pi_\theta(t)$};

    \nextgroupplot[xlabel={$f$}, xmin=-6,xmax=6, ymin=-0.3,ymax=1.15]
    \addplot[niceblue, thick, domain=-6:6, samples=200] {sin(deg(3.14159*x))/(3.14159*x)};
    \node at (axis cs: 0,1.35) {\small $\theta\sinc(\pi f\theta)$};

    \nextgroupplot[xlabel={$t$}, xmin=-6,xmax=6, ymin=-0.3,ymax=1.15]
    \addplot[red, thick, domain=-6:6, samples=200] {sin(deg(3.14159*x))/(3.14159*x)};
    \node at (axis cs: 0,1.35) {\small $\theta\sinc(\pi t\theta)$};

    \nextgroupplot[xlabel={$f$}, xmin=-2.2,xmax=2.2, ymin=-0.2,ymax=1.5]
    \addplot[red, thick] coordinates {(-2,0)(-0.5,0)(-0.5,1)(0.5,1)(0.5,0)(2,0)};
    \node at (axis cs: 0,1.7) {\small $\Pi_\theta(f)$};
    \end{groupplot}
    \node at ($(dualgrid c1r1.south)!0.5!(dualgrid c1r2.north)$)
        {\LARGE $\overunderset{\mathcal F}{\mathcal F^{-1}}{\rightleftarrows}$};
\end{tikzpicture}
